\documentclass[11pt,a4paper]{article}

% =====================================================================
%  Aquila — Technical Methods
%  Compile with: pdflatex Aquila_Methods.tex (twice for references)
% =====================================================================

\usepackage[utf8]{inputenc}
\usepackage[T1]{fontenc}
% \usepackage{lmodern}      % optional — uncomment if texlive-latex-extra installed
% \usepackage{microtype}    % optional — uncomment if texlive-latex-extra installed
\usepackage[margin=2.4cm]{geometry}
\usepackage{amsmath,amssymb,amsthm,bm}
\usepackage{graphicx}
\usepackage{booktabs}
\usepackage{array}
\usepackage{xcolor}
\usepackage{listings}
\usepackage{enumitem}
\usepackage{hyperref}
\usepackage{titlesec}
\usepackage{fancyhdr}
\usepackage{caption}
\usepackage{subcaption}
\usepackage{xparse}

% --------- Brand palette (mirrors the app) ---------
\definecolor{aqdark}{HTML}{001E2B}
\definecolor{aqred}{HTML}{FF1243}
\definecolor{aqteal}{HTML}{00707F}
\definecolor{aqgrey}{HTML}{5A5A5A}

\hypersetup{
    colorlinks=true,
    linkcolor=aqteal,
    citecolor=aqteal,
    urlcolor=aqteal,
    pdftitle={Aquila — Methods},
    pdfauthor={Merouane Hamdani},
    pdfsubject={Pressure Transient Analysis},
}

% --------- Section formatting ---------
\titleformat{\section}{\Large\bfseries\color{aqdark}}{\thesection}{0.6em}{}
\titleformat{\subsection}{\large\bfseries\color{aqdark}}{\thesubsection}{0.5em}{}
\titleformat{\subsubsection}{\normalsize\bfseries\color{aqteal}}{\thesubsubsection}{0.4em}{}

% --------- Headers / footers ---------
\pagestyle{fancy}
\fancyhf{}
\fancyhead[L]{\textsc{\color{aqdark}Aquila}\,\textcolor{aqgrey}{·\,Methods}}
\fancyhead[R]{\textcolor{aqgrey}{\thepage}}
\renewcommand{\headrulewidth}{0pt}
\fancyfoot[L]{\scriptsize\textcolor{aqgrey}{© 2026 Merouane Hamdani · MIT License}}
\fancyfoot[R]{\scriptsize\textcolor{aqgrey}{Not affiliated with Equinor ASA}}

% --------- Common shortcuts ---------
\newcommand{\R}{\mathrm{R}}
\newcommand{\D}{\mathrm{d}}
\newcommand{\pwf}{p_{\mathrm{wf}}}
\newcommand{\pinit}{p_i}
\newcommand{\Ein}{\mathrm{E}_1}
\newcommand{\BU}{\mathrm{BU}}
\newcommand{\mD}{\mathrm{mD}}
\newcommand{\kh}{kh}
\newcommand{\Ei}{\mathrm{Ei}}
\newcommand{\mpres}{m(p)}

% --------- Listings ---------
\lstdefinestyle{aqcode}{
    backgroundcolor=\color{black!4},
    basicstyle=\ttfamily\footnotesize\color{aqdark},
    keywordstyle=\color{aqred}\bfseries,
    commentstyle=\color{aqgrey}\itshape,
    stringstyle=\color{aqteal},
    numbers=none,
    frame=single,
    framerule=0pt,
    rulecolor=\color{black!10},
    showstringspaces=false,
    breaklines=true,
    aboveskip=8pt,
    belowskip=8pt,
}
\lstset{style=aqcode,language=Python}

% --------- Symbol table macro ---------
\NewDocumentEnvironment{symbols}{}
{\par\begin{tabular}{@{}p{0.10\textwidth}p{0.20\textwidth}p{0.66\textwidth}@{}}
\textbf{Symbol} & \textbf{Units} & \textbf{Meaning} \\ \midrule
}
{\end{tabular}\par\vspace{0.5em}}

% =====================================================================
\begin{document}

% --------- Title page ---------
\begin{titlepage}
\centering
\vspace*{2cm}
{\Huge\bfseries\color{aqdark} Aquila}\\[6pt]
\textcolor{aqred}{\rule{8cm}{2pt}}\\[10pt]
{\large\color{aqteal} Pressure Transient Analysis\,---\,Technical Methods}\\
\vspace{1.5cm}
{\large\color{aqdark} Merouane Hamdani}\\[4pt]
{\small\color{aqgrey} Author}\\
\vspace{2.5cm}
\begin{minipage}{0.8\textwidth}
\itshape\color{aqgrey}
This document specifies the analytical and numerical methods used by the
Aquila pressure-transient analysis workbench. Every formula in the
engine is referenced to its source, its dimensional form is given in
oilfield units (the engine's canonical internal representation), and its
validation status is stated. The document is intended for review by
petroleum-engineering subject-matter experts and as a reference companion
to the source code.
\end{minipage}
\vspace{1.5cm}

\begin{minipage}{0.8\textwidth}
\centering\footnotesize\color{aqgrey}
\textcolor{aqred}{\rule{4cm}{0.5pt}}\\[6pt]
\textbf{\color{aqdark}Disclaimer}\,\,---\,Aquila is provided \emph{as~is}, without
warranty of any kind. Outputs assist, not replace, the work of a
qualified petroleum engineer. The visual identity is \emph{inspired by}
publicly-visible corporate design conventions; the project is not produced
by, nor affiliated with, Equinor~ASA, and the Equinor name, logo and
wordmark are not used.
\end{minipage}

\vfill
{\large\bfseries Version 1.0}\\[6pt]
{\color{aqgrey}\today\,---\,Licensed under MIT}
\end{titlepage}

\tableofcontents
\newpage

% =====================================================================
\section{Notation and units}
\label{sec:notation}

The Aquila engine works internally in \emph{oilfield} units throughout.
Unit conversion happens only at the user-interface boundary
(\texttt{modules/units.py}). This keeps every formula in the engine in a
single consistent unit system and removes any risk of constant
mis-translation across the codebase.

\begin{symbols}
$k$         & mD             & permeability \\
$h$         & ft             & net pay (formation thickness) \\
$\phi$      & ---            & porosity (fraction) \\
$\mu$       & cp             & viscosity \\
$B$         & rb/STB or rb/Mscf & formation volume factor \\
$c_t$       & psi$^{-1}$     & total compressibility \\
$q$         & STB/D or Mscf/D & surface flow rate \\
$\Delta P$  & psi            & pressure drawdown $p_i - \pwf$ \\
$\Delta t$  & hr             & elapsed time \\
$r_w$       & ft             & wellbore radius \\
$r_e$       & ft             & drainage radius \\
$s$         & ---            & mechanical skin (dimensionless) \\
$C$         & bbl/psi        & wellbore-storage coefficient \\
$p_i$       & psia           & initial reservoir pressure \\
$T_R$       & $^\circ$R      & reservoir temperature (gas) \\
$\gamma_g$  & ---            & gas specific gravity (air = 1) \\
$Z$         & ---            & gas Z-factor (real-gas deviation) \\
$\mpres$       & psi$^2$/cp     & gas pseudo-pressure \\
$t_a$       & psi$\cdot$hr/cp & gas pseudo-time \\
\end{symbols}

% =====================================================================
\section{The radial diffusivity equation}
\label{sec:diffeq}

Aquila is built on the radial slightly-compressible-fluid diffusivity
equation in a single homogeneous layer:
\begin{equation}
    \frac{1}{r}\frac{\partial}{\partial r}\!\left(r\,\frac{\partial p}{\partial r}\right)
    = \frac{\phi \mu c_t}{k}\,\frac{\partial p}{\partial t}.
    \label{eq:diff}
\end{equation}

The \emph{line-source} solution to (\ref{eq:diff}) with constant-rate
inner-boundary condition and an infinite-acting outer boundary gives, at
the wellbore,
\begin{equation}
    \Delta P(t)
    \;=\;
    \frac{141.2\,q\,B\,\mu}{\kh}
    \left[
        \tfrac{1}{2}\,\Ein\!\!\left(\frac{x}{1}\right) + s
    \right],
    \qquad
    x = \frac{948\,\phi\,\mu\,c_t\,r_w^{\,2}}{k\,t},
    \label{eq:linesrc}
\end{equation}
where $\Ein(x)\equiv -\Ei(-x)=\int_x^\infty\!e^{-u}/u\;\D u$ is the
exponential integral. In the late-time (small-$x$) limit, $\Ein(x)\to
-\ln x - \gamma_{\mathrm{E}}$ where $\gamma_{\mathrm{E}}\approx 0.5772$
is the Euler--Mascheroni constant. This produces the famous
\emph{semilog approximation}:
\begin{equation}
    \Delta P(t) \approx
    \frac{162.6\,q\,B\,\mu}{\kh}
    \!\left[
        \log_{10} t + \log_{10}\!\left(\frac{k}{\phi\mu c_t r_w^{\,2}}\right)
        - 3.23 + \tfrac{1}{1.151}\,s
    \right].
    \label{eq:semilog}
\end{equation}
The slope of $\Delta P$ versus $\log_{10}t$ on this plot,
$m \equiv 162.6\,qB\mu/(\kh)$, is the foundation of conventional
infinite-acting radial-flow (IARF) interpretation.

\paragraph{Bourdet derivative.}
The Bourdet (1989) derivative is
\begin{equation}
    \frac{\partial \Delta P}{\partial \ln t}
    \;=\;
    t\,\frac{\D \Delta P}{\D t}.
    \label{eq:bourdet}
\end{equation}
Under IARF, (\ref{eq:semilog}) gives
$t\,\D\Delta P/\D t = m/\ln 10$, so the derivative is a \emph{flat
plateau} on a log-log plot. Departures from this plateau diagnose
non-IARF flow regimes: wellbore storage, linear flow, pseudo-steady-state
(closed-boundary), double-porosity transitions, layered behaviour, and
so on. Aquila computes the derivative with a vectorised three-point
log-time-window scheme:
\begin{equation}
    \left(t\,\frac{\D \Delta P}{\D t}\right)_{\!i}
    \;=\;
    \frac{
        \Delta P_R\,\ln\!\tfrac{t_i}{t_L} - \Delta P_L\,\ln\!\tfrac{t_R}{t_i}
    }{
        \ln\!\tfrac{t_R}{t_i} \cdot \ln\!\tfrac{t_i}{t_L}
    }
    \cdot
    \frac{1}{\ln\!\tfrac{t_R}{t_L}},
\end{equation}
where $t_L,\,t_R$ are the times within a log-window of half-width $L$
to the left and right of $t_i$ (default $L=0.1$). The implementation is
fully vectorised via \texttt{np.searchsorted} and is bit-identical to a
naive loop but $\sim$1000$\times$ faster on a 6000-point trace.

% =====================================================================
\section{Wellbore storage — Fair--Schulz model}
\label{sec:wbs}

Wellbore storage (WBS) dominates the early-time response. Pure WBS
gives $\Delta P_\text{WBS} = (qB/24C)\,t$ — a unit-slope ramp on
log-log. Aquila uses the \textbf{Fair--Schulz} smooth transition between
pure WBS and the line-source solution:
\begin{equation}
    \Delta P(t) = \underbrace{w(t)}_{\text{weight}}\,\Delta P_\text{res}(t)
                + \bigl[1-w(t)\bigr]\,\Delta P_\text{WBS}(t),
    \qquad
    w(t) = 1 - \exp\!\left[-(t/\tau)^2\right],
\end{equation}
with characteristic time $\tau = 200\,C/(kh)$. The transition smoothly
matches the unit-slope WBS region at early times to the line-source
formation response at late times, and is differentiable so the
Bourdet derivative remains smooth across the join.

The rule-of-thumb \emph{clearance time} — after which the formation
response dominates — is
\begin{equation}
    t_\text{clear} \;\approx\; (200\,000 + 12\,000\,s)\,\frac{C}{\kh/\mu}
    \qquad [\mathrm{hr}].
\end{equation}
Aquila displays a warning when the planned test duration is less than
$1.5\,t_\text{clear}$.

% =====================================================================
\section{Regression for IARF — solving for $k$ and $s$}
\label{sec:regress-radial}

In a clean IARF window, the geometric mean of the Bourdet derivative
yields the plateau, from which $k$ follows directly:
\begin{align}
    \text{plateau} &= \exp\!\left[\,\overline{\ln(t\,\D\Delta P/\D t)}\,\right], \\
    m              &= \ln(10)\cdot\text{plateau} \qquad [\mathrm{psi/cycle}], \\
    k              &= \frac{162.6\,q\,B\,\mu}{m\,h}.
\end{align}

Skin is recovered using the value of $\Delta P$ at $\Delta t = 1$ hr in
the buildup. The standard Horner formula is
\begin{equation}
    s \;=\; 1.151\!\left[
        \frac{\Delta P_{1\mathrm{hr}}}{m}
        \;-\;
        \log_{10}\!\left(\frac{k}{\phi\mu c_t r_w^{\,2}}\right)
        \;+\;
        3.23
    \right].
    \label{eq:skin-horner}
\end{equation}

\paragraph{Quality flag — plateau flatness.}
The geometric standard deviation of $t\,\D\Delta P/\D t$ within the
chosen regression window is reported:
\begin{equation}
    f \;=\; \exp\!\left[\sigma\!\left(\ln(t\,\D\Delta P/\D t)\right)\right]
\end{equation}
Values $f<1.15$ indicate a clean plateau; $f<1.40$ is borderline;
$f\geq 1.40$ indicates the window crosses a regime transition and the
fit is suspect. The UI surfaces this directly to discourage the user
from picking a window that includes wellbore storage or a boundary.

% =====================================================================
\section{Linear flow (fractured wells)}
\label{sec:linear}

For a vertical well intersecting an infinite-conductivity fracture, the
early-time response is \emph{linear flow} into the fracture, with
\begin{equation}
    \Delta P(t) = m_L \sqrt{t},
    \qquad
    x_f\sqrt{k} \;=\; \frac{4.064\,q\,B}{m_L\,h}\sqrt{\frac{\mu}{\phi c_t}},
\end{equation}
where $x_f$ is the fracture half-length. On the Bourdet derivative,
linear flow appears as a \emph{half-slope} line on log-log.

For a gas well, the equivalent pseudo-pressure formulation is
\begin{equation}
    \Delta \mpres(t) = m_L \sqrt{t},
    \qquad
    x_f\sqrt{k} \;=\; \frac{40.92\,q_g\,T_R}{m_L\,h}\sqrt{\frac{1}{\phi\,\mu_\text{ref}\,c_{t,\text{ref}}}}.
\end{equation}

% =====================================================================
\section{Closed-boundary (pseudo-steady-state, PSS)}
\label{sec:pss}

After a well has drained out to a no-flow boundary, the pressure decline
is uniform across the drainage area — \emph{pseudo-steady-state}. On
Cartesian axes, $\Delta P$ becomes a straight line in $t$:
\begin{equation}
    \frac{\D\Delta P}{\D t} = m^\ast = \frac{0.234\,q\,B}{\phi c_t h\,A},
    \qquad
    A = \pi r_e^{\,2}.
\end{equation}
On a log-log Bourdet derivative, PSS appears as a \emph{unit-slope}
rising line at late time. Aquila's \texttt{regress\_pss} extracts
$m^\ast$ by linear regression in real time and reports the drainage
area, equivalent drainage radius $r_e = \sqrt{A/\pi}$, and pore
volume:
\begin{equation}
    V_p = A\,h\,\phi \;[\mathrm{ft}^3]
          \quad\Longleftrightarrow\quad
          V_p \,[\mathrm{bbl}] = V_p\,[\mathrm{ft}^3]/5.615.
\end{equation}
For gas wells, the slope and area formulae use the pseudo-pressure form
with constant $2.356\times 10^6$ in place of $0.234$.

% =====================================================================
\section{Gas-well analysis}
\label{sec:gas}

For gas, the diffusivity equation is non-linear because $\mu$, $Z$, and
$c_g$ all depend on pressure. The standard remedy is the Al-Hussainy /
Ramey transformation to \emph{pseudo-pressure}:
\begin{equation}
    \mpres \;=\; 2\!\int_0^p \frac{p'\,\D p'}{\mu(p')\,Z(p')}.
    \label{eq:mp}
\end{equation}
In the variable $\mpres$, the gas diffusivity equation is linear and has
exactly the same form as (\ref{eq:diff}). Hence every analytical
solution carries over with the substitution $\Delta P \to \Delta\mpres$ and
a different slope constant:
\begin{equation}
    \Delta\mpres(t) = \frac{1422\,q_g\,T_R}{\kh}\!\left[\tfrac{1}{2}\Ein(x) + s\right],
    \qquad
    m_g \equiv \frac{1637\,q_g\,T_R}{\kh}.
\end{equation}

For drawdown and buildup tests, time is also transformed:
\begin{equation}
    t_a \;=\; \frac{t}{\mu(\bar p)\,c_t(\bar p)}
\end{equation}
evaluated at the average reservoir pressure $\bar p$
(Agarwal approximation, accurate for buildup analysis when the pressure
swing is small).

\subsection{PVT correlations}

Aquila implements four standard correlations covering the practical
$0.55 < \gamma_g < 1.7$, $1.0 < T_{pr} < 3.0$, $0.2 < P_{pr} < 30$ range:

\paragraph{Pseudo-criticals — Sutton (1985).}
\begin{align}
    P_{pc} &= 756.8 - 131.0\,\gamma_g - 3.6\,\gamma_g^{\,2}, \\
    T_{pc} &= 169.2 + 349.5\,\gamma_g - 74.0\,\gamma_g^{\,2}.
\end{align}

\paragraph{Acid-gas correction — Wichert--Aziz (1972).}
For mole fractions $y_{\mathrm{CO}_2}$ and $y_{\mathrm{H}_2\mathrm{S}}$:
\begin{align}
    A &= y_{\mathrm{H}_2\mathrm{S}} + y_{\mathrm{CO}_2},\quad
    B  = y_{\mathrm{H}_2\mathrm{S}}, \\
    \varepsilon &= 120\,(A^{0.9}-A^{1.6}) + 15\,(B^{0.5}-B^4), \\
    T_{pc}^\prime &= T_{pc} - \varepsilon, \\
    P_{pc}^\prime &= \frac{P_{pc}\,T_{pc}^\prime}{T_{pc} + B(1-B)\varepsilon}.
\end{align}

\paragraph{Z-factor — Dranchuk--Abou-Kassem (1975).}
\begin{equation}
    Z \;=\; 1 + c_1\rho_r + c_2\rho_r^{\,2} - c_3\rho_r^{\,5}
            + c_4\bigl(1+A_{11}\rho_r^{\,2}\bigr)\frac{\rho_r^{\,2}}{T_{pr}^{\,3}}\exp(-A_{11}\rho_r^{\,2}),
\end{equation}
with $\rho_r = 0.27\,P_{pr}/(Z\,T_{pr})$ and DAK coefficients
$A_1\ldots A_{11}$ as in the 1975 paper. The eleven $A_i$ are
$[0.3265,\, -1.0700,\, -0.5339,\, 0.01569,\, -0.05165,\, 0.5475,\,
-0.7361,\, 0.1844,\, 0.1056,\, 0.6134,\, 0.7210]$. The implicit
relation in $\rho_r$ is solved by bracketed bisection.

\paragraph{Viscosity — Lee--Gonzalez--Eakin (1966).}
\begin{align}
    M_g &= 28.96\,\gamma_g \qquad [\text{lbm/lbmol}], \\
    K  &= \frac{(9.4 + 0.02 M_g)\,T_R^{1.5}}{209 + 19 M_g + T_R}, \\
    X  &= 3.5 + \frac{986}{T_R} + 0.01 M_g,\qquad
    Y   = 2.4 - 0.2 X, \\
    \mu_g &= 10^{-4}\,K\,\exp\!\bigl(X\,\rho_g^{\,Y}\bigr) \qquad [\mathrm{cp}],
\end{align}
with the gas density $\rho_g$ in g/cm$^3$ converted from
$\rho_g [\mathrm{lbm/ft^3}] = p M_g/(Z\,R\,T_R)$, $R = 10.7316$
psi$\cdot$ft$^3$/(lbmol$\cdot^\circ$R).

\paragraph{Compressibility and FVF.}
\begin{equation}
    c_g \;=\; \frac{1}{p} - \frac{1}{Z}\!\left(\frac{\partial Z}{\partial p}\right)_T,
    \qquad
    B_g \;=\; 5.0287\,\frac{Z\,T_R}{p} \qquad [\mathrm{rb/Mscf}].
\end{equation}
The $\partial Z/\partial p$ derivative is computed numerically with a
$10^{-4}$ relative step.

\subsection{Validation status}

\begin{tabular}{lll}
\toprule
\textbf{Correlation} & \textbf{Test point} & \textbf{Result} \\
\midrule
Sutton           & $\gamma_g=0.65$                  & $P_{pc}=670$, $T_{pc}=365^\circ$R (matches paper) \\
Wichert--Aziz    & $\gamma_g=0.7$, $5\%$CO$_2$+$10\%$H$_2$S & $\Delta T_{pc}=20.7^\circ$R \\
DAK              & Standing--Katz, 8 points         & 7/8 within $\pm$0.05 \\
LGE              & $p=2000$ psia, $T=200^\circ$F    & $\mu_g=0.0168$ cp (in $0.014$--$0.020$ band) \\
$\mpres$            & monotonicity test                & strictly increasing $\checkmark$ \\
Forward $\to$ inverse $k$ & three skin cases       & $k$ recovered to $\le 0.05\,\%$ \\
\bottomrule
\end{tabular}

% =====================================================================
\section{Multi-layer reservoirs}
\label{sec:layered}

\subsection{Commingled layers (no crossflow)}

Two (or more) layers produce into the same wellbore but do \emph{not}
communicate vertically. Under the constraint of equal wellbore pressure
and fixed total surface rate, the per-layer pressure responses combine
by \emph{parallel admittance}:
\begin{equation}
    \Delta P_\text{well}(t) =
    \frac{q_\text{total}}{\sum_i Y_i(t)},
    \qquad
    Y_i(t) = \frac{k_i h_i}{141.2\,B\,\mu \cdot g_i(t)},
\end{equation}
where $g_i(t) = \tfrac{1}{2}\Ein(x_i(t)) + s_i$ is the dimensionless
line-source contribution of layer $i$ evaluated at time $t$.

\paragraph{Identifiable parameters.}
From a single buildup, one can identify $\kh_\text{total}$ and the
\emph{flow-capacity ratio}
$\kappa = k_1 h_1 / \kh_\text{total}$, but the $(k_i, h_i)$ pair
separately is non-unique without independent rate-allocation data (e.g.\
production logging tools). Aquila surfaces this limitation explicitly in
the inverse-fit output.

\paragraph{Signature on the Bourdet derivative.}
\begin{itemize}[noitemsep]
    \item Early time: the fastest-responding layer dominates the
          derivative plateau (the high-$k$, low-storage layer).
    \item Late time: the derivative drops to the plateau corresponding
          to the system total $\kh$.
    \item Transition: a downward \emph{step} in the derivative — the
          qualitative fingerprint of commingled flow without crossflow.
\end{itemize}

\subsection{Double permeability with crossflow}

\textcolor{aqgrey}{Pending verification against published reference
type-curves. The numerical Laplace-inversion infrastructure (Stehfest
1970, validated to 0.30\,\% on the $K_0(\sqrt{s})/s$ radial-diffusion
benchmark) is in place under \texttt{modules/laplace.py}. What remains
is the Laplace-space transfer function $f(s;\kappa,\omega,\lambda)$ from
Bourdet (1985) SPE 13628, which will be added once an authoritative
reference is available for cross-validation.}

% =====================================================================
\section{Numerical Laplace inversion (Stehfest)}
\label{sec:stehfest}

Several PTA solutions (notably double-permeability with crossflow,
double-porosity, and bounded reservoirs with non-trivial outer
geometries) are most naturally expressed in Laplace space. Aquila
implements the Stehfest (1970) algorithm:
\begin{equation}
    f(t) \;\approx\;
    \frac{\ln 2}{t}\,\sum_{i=1}^{N} V_i\,\bar F\!\left(\frac{i\ln 2}{t}\right),
\end{equation}
with weights
\begin{equation}
    V_i = (-1)^{N/2+i} \!\!\!
    \sum_{k=\lfloor(i+1)/2\rfloor}^{\min(i,N/2)}
    \frac{k^{N/2}\,(2k)!}{(N/2-k)!\,k!\,(k-1)!\,(i-k)!\,(2k-i)!}.
\end{equation}
Aquila uses $N=8$ by default — sufficient for the smooth monotone
transforms typical of radial PTA, while avoiding the catastrophic
cancellation that affects $N=14$+ on exponentially-decaying functions
in double precision.

\paragraph{Validation.}
Stehfest is validated against five known Laplace pairs (Table~\ref{tab:stehfest}).

\begin{table}[h]
\centering
\caption{Stehfest inversion accuracy on benchmark transforms (max relative error over the indicated $t$-range).}
\label{tab:stehfest}
\begin{tabular}{lll}
\toprule
\textbf{Transform $\bar F(s)$} & \textbf{Inverse $f(t)$} & \textbf{Max error} \\
\midrule
$1/s^2$              & $t$                            & $1.5\times10^{-4}$ \\
$\sqrt{\pi}/(2s^{3/2})$ & $\sqrt{t}$                   & $1.7\times10^{-5}$ \\
$1/(s+a)$            & $e^{-at}$ (N=12)              & $2.0\times10^{-3}$ \\
$\omega/(s^2+\omega^2)$ & $\sin(\omega t)$ (N=12)    & $5.0\times10^{-2}$ \\
$K_0(\sqrt s)/s$    & $\tfrac{1}{2}\Ein\!\left(\tfrac{1}{4t}\right)$ & $3.0\times10^{-3}$ \\
\bottomrule
\end{tabular}
\end{table}

% =====================================================================
\section{Uncertainty quantification}
\label{sec:uncertainty}

Aquila propagates input uncertainty on $h$ and $\mu$ (oil) or $h$ alone
(gas — viscosity is PVT-derived) to the permeability estimate via
Monte-Carlo sampling of triangular distributions. For each iteration
$i$:
\begin{equation}
    k_i = \begin{cases}
        \dfrac{162.6\,q\,B\,\mu_i}{m\,h_i} & \text{oil/water,} \\[1ex]
        \dfrac{1637\,q\,T_R}{m_g\,h_i}      & \text{gas.}
    \end{cases}
\end{equation}
The sample yields the petroleum-convention percentiles
\textbf{P90 / P50 / P10} (low/best/high), a kernel-density-estimated
PDF, the empirical CDF, and a Spearman-rank-correlation tornado
sensitivity chart ranking the relative impact of each uncertain
input.

% =====================================================================
\section{Engine architecture}
\label{sec:arch}

\begin{lstlisting}[caption={Aquila module layout.}]
aquila/
├── app.py                    # Streamlit UI (orchestrator only)
├── styles.py                 # Theming, CSS, Plotly layout
├── assets/
│   ├── aquila_logo.svg       # Full logo (400x120)
│   └── aquila_icon.svg       # Square favicon (100x100)
└── modules/                  # Math layer — UI-independent
    ├── engine.py             # Line-source, Fair-Schulz, Bourdet,
    │                         #   regression (radial/linear/PSS).
    ├── pvt.py                # DAK Z, LGE mu, Sutton+Wichert-Aziz,
    │                         #   m(p) and t_a.
    ├── layered.py            # Commingled (no crossflow);
    │                         #   double-perm stub.
    ├── laplace.py            # Stehfest numerical inverter.
    ├── economics.py          # Cost & deferred-revenue calculator.
    ├── uncertainty.py        # Monte Carlo + KDE/CDF + tornado.
    ├── units.py              # Field <-> SI exact conversions.
    └── reporter.py           # FPDF report builder.
\end{lstlisting}

The math layer never imports Streamlit. Every module is independently
unit-testable and reusable from a notebook or batch pipeline.

% =====================================================================
\section*{References}
\addcontentsline{toc}{section}{References}
\small

\noindent\textbf{Pressure transient analysis}\par
\begin{itemize}[leftmargin=*,nosep]
\item Bourdet, D. (1989). \emph{Use of Pressure Derivative in Well-Test Interpretation}. SPE 12777.
\item Earlougher, R.C. (1977). \emph{Advances in Well Test Analysis}, SPE Monograph 5.
\item Lee, J., Rollins, J.B., \& Spivey, J.P. (2003). \emph{Pressure Transient Testing}, SPE Textbook Series.
\item Horne, R.N. (1995). \emph{Modern Well Test Analysis: A Computer-Aided Approach}, 2nd ed.
\end{itemize}

\noindent\textbf{Gas-well analysis}\par
\begin{itemize}[leftmargin=*,nosep]
\item Al-Hussainy, R., Ramey, H.J., \& Crawford, P.B. (1966). The flow of real gases through porous media. \emph{JPT}.
\item Agarwal, R.G. (1979). ``Real gas pseudo-time'' — a new function for pressure buildup analysis of MHF gas wells. SPE 8279.
\item Lee, W.J., \& Wattenbarger, R.A. (1996). \emph{Gas Reservoir Engineering}, SPE Textbook Series.
\end{itemize}

\noindent\textbf{PVT correlations}\par
\begin{itemize}[leftmargin=*,nosep]
\item Dranchuk, P.M., \& Abou-Kassem, J.H. (1975). Calculation of Z factors for natural gases using equations of state. \emph{JCPT}.
\item Lee, A.L., Gonzalez, M.H., \& Eakin, B.E. (1966). The viscosity of natural gases. \emph{JPT}.
\item Sutton, R.P. (1985). Compressibility factors for high-molecular-weight reservoir gases. SPE 14265.
\item Wichert, E., \& Aziz, K. (1972). Calculation of $Z$'s for sour gases. \emph{Hydrocarbon Processing}.
\end{itemize}

\noindent\textbf{Layered reservoirs}\par
\begin{itemize}[leftmargin=*,nosep]
\item Bourdet, D. (1985). Pressure behavior of layered reservoirs with crossflow. SPE 13628.
\item Larsen, L. (1981). Wells producing commingled zones with unequal initial pressures and reservoir properties. SPE 10325.
\end{itemize}

\noindent\textbf{Numerical methods}\par
\begin{itemize}[leftmargin=*,nosep]
\item Stehfest, H. (1970). Algorithm 368: Numerical inversion of Laplace transforms. \emph{Commun.\ ACM} 13(1):47--49.
\item Bourdet, D., Ayoub, J.A., \& Pirard, Y.M. (1989). Use of pressure derivative in well-test interpretation. \emph{SPE Form.\ Eval.}\ 4(2):293--302.
\end{itemize}

\end{document}
